\section{数值方法}\label{sec:numerics}

第~\ref{sec:theory} 节把成像问题化归为一条一维二阶常微分方程
\eqref{eq:geodesic}。本节说明模拟器如何在 GPU 上以每像素一次的代价把它解到
可验证的精度：如何从相机参数构造轨道平面的初值（\S\ref{sec:plane}）、
用什么格式推进（\S\ref{sec:rk4}）、步长如何自适应（\S\ref{sec:stepsize}）、
何时终止（\S\ref{sec:termination}），以及沿光线如何累积盘辐射
（\S\ref{sec:rtnum}）。最后 \S\ref{sec:residual} 给出守恒残差的定义与
后文验证实验的接口。

\subsection{从像素到轨道平面}\label{sec:plane}

\paragraph{观测者标架。}
相机是施瓦西时空中的一个静态观测者，其坐标位置记为 $\vec r_0$，
$r_0=|\vec r_0|$，径向单位向量 $\hat e_1=\vec r_0/r_0$。
按式~\eqref{eq:tetrad}，该观测者的正交归一四元标架把坐标分量与局域
「实验室」分量联系起来。着色器为每个像素构造一个局域单位方向
\begin{equation}
n^{\hat a}_{\rm local}
=\mathcal N\!\left(\hat e_{\hat a}\cdot
\bigl(\hat x\,n_x\tan\tfrac{\theta_{\rm fov}}{2}\,A,\;
\hat y\,n_y\tan\tfrac{\theta_{\rm fov}}{2},\;
\hat z\bigr)\right),
\label{eq:nlocal}
\end{equation}
其中 $n_x,n_y\in[-1,1]$ 是该像素的归一化设备坐标（NDC），$A$ 为视口宽高比，
$\theta_{\rm fov}$ 为垂直视场角，$\mathcal N$ 表示归一化。像素抖动
（\S\ref{sec:jitter}）在 $n_{x,y}$ 上加一个低于一像素的随机偏置，
使多帧累积后的阴影边缘不出现阶梯状锯齿。

\paragraph{轨道平面。}
测地线是平面的：由 $\vec r_0$ 与 $n_{\rm local}$ 张成的平面包含整条光线。
其法向量为
$\vec c_N=\vec r_0\times n_{\rm local}$，平面内正交基取
\begin{equation}
\hat E_1=\frac{\vec r_0}{r_0},\qquad
\hat E_2=\frac{\vec c_N\times \hat E_1}{|\vec c_N\times\hat E_1|},
\qquad \text{若 } n_{\rm local}\cdot\hat E_2<0 \text{ 则 } \hat E_2\to-\hat E_2 .
\label{eq:plane_basis}
\end{equation}
最后一步的符号修正是必要的：$\hat E_2$ 的方向由叉乘确定，
在 $\vec c_N$ 与 $\hat E_1$ 接近共线时会由数值噪声决定，而积分方向
$\varphi$ 递增必须对应光线的前进方向。有了相位约定后，光线的坐标轨迹
就是一条平面曲线
\begin{equation}
\vec p(\varphi)=\frac{\cos\varphi\,\hat E_1+\sin\varphi\,\hat E_2}{u(\varphi)},
\qquad u\equiv\frac{1}{r},
\label{eq:plane_traj}
\end{equation}
于是 $\vec p(0)=\vec r_0$ 自动成立，整条测地线由标量函数 $u(\varphi)$ 完全决定。
这正是把二维光线追踪压缩成一维积分的全部理由。

\paragraph{初值：从局域方向到坐标斜率。}
设 $a=n_{\rm local}\cdot\hat E_1=\cos\psi$，则
$\sin\psi=\sqrt{1-a^2}$。局域方向中径向分量与切向分量之比为
$a:\sin\psi$。径向固有长度为 $\dd \ell_r=\dd r/\sqrt{1-\rs/r}$，
切向固有长度为 $r\,\dd\varphi$，二者在局域标架中是同一组分量，
因此在坐标中
\begin{equation}
\frac{\dd r}{\dd\varphi}
=r\sqrt{1-\frac{\rs}{r}}\;\frac{a}{\sin\psi}
=r_0\sqrt{f_0}\;\frac{a}{\sqrt{1-a^2}},
\qquad f_0\equiv 1-\frac{\rs}{r_0},
\label{eq:drdphi}
\end{equation}
由 $u=1/r$ 得 $\dd u/\dd\varphi=-u^2\,\dd r/\dd\varphi$，代入 $u_0=1/r_0$：
\begin{equation}
u(0)=u_0=\frac{1}{r_0},\qquad
\left.\frac{\dd u}{\dd\varphi}\right|_{\varphi=0}
=-u_0\,\frac{a}{\sqrt{1-a^2}}\,\sqrt{f_0}.
\label{eq:ic}
\end{equation}
式~\eqref{eq:ic} 与着色器中
\texttt{du = -u0 * (a/btan) * fObs} 逐字符对应。

\paragraph{冲击参数作为第一积分。}
式~\eqref{eq:EL} 的第一积分给出了这条光线的冲击参数
\begin{equation}
b=\Bigl[\bigl(\tfrac{\dd u}{\dd\varphi}\bigr)^{2}
+u^{2}\bigl(1-2\M u\bigr)\Bigr]^{-1/2},
\label{eq:bfromic}
\end{equation}
它是一个沿光线严格守恒的量，在着色器中即为 \texttt{G\_B}。
$b$ 同时是后续多个模块的关键输入：阴影判据 $b\lessgtr \bc$、
盘辐射的轴向多普勒投影（式~\eqref{eq:baxis}）以及验证实验 V1 的被测量。
把式~\eqref{eq:ic} 代入式~\eqref{eq:bfromic}，
\begin{equation}
b=\frac{r_0\sin\psi}{\sqrt{f_0}},
\label{eq:b_sinpsi}
\end{equation}
即「视线与径向的夹角」到「冲击参数」的解析映射。
式~\eqref{eq:b_sinpsi} 在 V4 中被用作独立校验：由视线几何反解出的 $b$
与积分器首步的第一积分必须一致到浮点精度（表~\ref{tab:v4a}）。

\paragraph{退化情形。}
当 $\sin\psi<10^{-5}$ 时平面 $\vec r_0\wedge n_{\rm local}$ 退化，
式~\eqref{eq:plane_basis} 中的叉乘失去意义。这类光线的轨迹是一条
过原点的直线：若 $a<0$（径向向内）必然落入视界，输出纯黑；
若 $a>0$（径向向外）永不偏折，直接输出背景星场。
两种情况都在着色器里由解析分支处理，既避免数值爆炸，
也不会在图像的正中心留下一个数值伪影。

\subsection{四阶 Runge--Kutta 格式}\label{sec:rk4}

把式~\eqref{eq:geodesic} 写成一阶系统
$\mathbf y=(u,w)^{\mathsf T}$，$w\equiv\dd u/\dd\varphi$：
\begin{equation}
\frac{\dd}{\dd\varphi}
\begin{pmatrix}u\\ w\end{pmatrix}
=\mathbf f(u,w)=\begin{pmatrix}w\\ -u+3\M u^{2}\end{pmatrix}.
\label{eq:system}
\end{equation}
右端是一个二次多项式，没有奇点、没有超越函数，非常适合在
着色器里高吞吐求值。模拟器采用经典四阶 Runge--Kutta（RK4）：
\begin{equation}
\begin{aligned}
\mathbf k_1&=\mathbf f(\mathbf y), &
\mathbf k_2&=\mathbf f(\mathbf y+\tfrac h2\mathbf k_1),\\
\mathbf k_3&=\mathbf f(\mathbf y+\tfrac h2\mathbf k_2), &
\mathbf k_4&=\mathbf f(\mathbf y+h\,\mathbf k_3),\\
\mathbf y_{n+1}&=\mathbf y_n+\frac h6\bigl(\mathbf k_1+2\mathbf k_2+2\mathbf k_3+\mathbf k_4\bigr).
\end{aligned}
\label{eq:rk4}
\end{equation}
RK4 的局部截断误差为 $\mathcal O(h^5)$、整体为 $\mathcal O(h^4)$。
选择 RK4 而不是更高阶格式有三个理由：\\
（一）每条光线平均只走 $10^2$--$10^3$ 步，四阶格式已经把误差压到
低于单精度浮点的舍入水平（\S\ref{sec:residual}，图~\ref{fig:05}）；
（二）RK4 每步 4 次右端求值、无历史状态，寄存器压力小，
是 GPU 上最不容易掉进分支发散与溢出陷阱的格式；
（三）$u$ 与 $w$ 各自的量级差异大（远场 $u\sim10^{-2}$、近视界 $u\sim0.5$），
高阶格式（如 Dormand--Prince）的自适应机制在此处收益有限，
却需要额外的误差估计与更多临时变量。

需要强调的是：积分变量是 $\varphi$ 而不是仿射参数 $\lambda$。以 $\varphi$ 为
自变量的代价是步长不再度量固有时间，但它带来了两个决定性好处：
（一）积分区间在平面内是单调的，不会因为 $r$ 的非单调行为在
近心点附近退化；（二）盘采样所需的水平坐标可以直接由
式~\eqref{eq:plane_traj} 给出，无需再积分一次位置方程。

\subsection{几何自适应步长}\label{sec:stepsize}

固定步长在施瓦西测地线上效率极低：远场几乎是一维谐振子
（$u''=-u$，周期 $2\pi$），可以用 $h\sim0.3$ 的大步推进；
而接近光子球时 $|w|$ 急剧变小、轨道弯曲极快，需要 $h\sim10^{-3}$。
模拟器采用「限制 $u$ 的相对变化率」的自适应律
\begin{equation}
h=\operatorname{clamp}\!\Bigl(\frac{tol\cdot u}{\max(|w|,10^{-7})},\;
0.002,\;h_{\max}(u)\Bigr)\times\sigma_h,
\qquad
h_{\max}(u)=0.35-0.30\,\mathrm{smoothstep}(0,0.30,u),
\label{eq:hstep}
\end{equation}
其中 $tol$ 是用户在界面上可调（默认 $0.03$）的相对容差，
$\sigma_h$ 是全局步长缩放（默认 $1$），$\mathrm{smoothstep}$ 为
$3x^2-2x^3$ 型三次平滑过渡。三条判据的作用分别是：

\begin{itemize}\itemsep2pt
\item \textbf{相对变化率} $h=tol\,u/|w|$：把单步内 $|\Delta u/u|$
  限制在 $tol$ 量级。远场 $|w|\ll u$，该式给出大步；近心点
  $|w|\to0$ 时该式发散，因此必须配合上界与下界使用。
\item \textbf{上界} $h_{\max}(u)$：$u\to0$（远场）时取 $0.35$\,rad，
  约为谐振子周期的 $5.6\%$，单步相位误差 $\sim10^{-5}$；
  $u\to0.30$ 以上（$r\lesssim3.3\M$，即光子球附近）平滑收窄到
  $0.05$\,rad。这个 $-0.30\,\mathrm{smoothstep}$ 的形状来自
  「误差随 $3\M u$ 增大而增大」的观察，用一个代价极低的表达式
  替代了完整的误差估计器。
\item \textbf{下界} $0.002$\,rad：保证在转折点（$w\approx0$）附近
  步长不会趋近于零而耗尽步数预算。这个下界也是本模拟器
  最主要的性能瓶颈来源，其代价在 \S\ref{sec:perf} 中用
  \texttt{diag\_taucap.py} 的对照实验定量给出。
\end{itemize}

由式~\eqref{eq:hstep} 可以立刻读出一个重要的等价性：把 $h$ 乘上
$\sigma_h$ 与把 $tol$ 除以 $\sigma_h$ 在 $h$ 未触界时**完全等价**。
因此在界面性能面板上 $tol$ 与 $\sigma_h$ 并非两个独立旋钮；
后文表~\ref{tab:perftol} 与图~\ref{fig:22} 用实验证实了这一点
（$\sigma_h=2$ 与 $tol\to tol/2$ 的耗时落在同一水平），
这也是我们把 $\sigma_h$ 默认值固定为 $1$ 的原因。

%% generated by tools/make_tables.py -- do not edit
\begin{table}[t]
\centering
\caption{光测地线积分器的全部数值判据。左列为源码中的字面写法，右列为物理含义；本表由 \texttt{tools/make\_tables.py} 直接匹配 \texttt{web/src/shaders.js} 的正则生成，任一常量被改动而未被本表捕获时编译脚本会立即报错。}
\label{tab:integrator}
\scriptsize
\begin{tabular}{p{0.34\linewidth}p{0.58\linewidth}}
\toprule
源码字面量 & 作用与取值 \\
\midrule
\texttt{for (int i = 0; i < N; i++)} & 编译期迭代硬上限（防止病态光线挂死 GPU） \\
\texttt{uMaxSteps} & 运行期步数上限 $n_{\max}$，用户可在 60--900 间调节 \\
\texttt{u > 0.5} & 捕获判据：$r<\rs$ 后不再有可能逃逸 \\
\texttt{u < uMin \&\& du < 0} & 逃逸判据：向外运动且已越过 $R_{\rm esc}$ \\
\texttt{uMin = 1/max(uEscapeR, 1.6 * r0)} & 逃逸半径至少为观测者半径的 1.6 倍，避免相机贴视界时的伪逃逸 \\
\texttt{hmax} & 步长上界：远场 $0.35$ rad，近视界收窄到 $0.05$ rad \\
\texttt{h = uTol * u / |du|} & 把 $|\Delta u/u|$ 限制在 $tol$ 量级（防止除零） \\
\texttt{clamp(h, 0.002, hmax)} & 步长下界固定为 $0.002$ rad：转折点附近 $|\dd u|\to0$ 时不无限细分 \\
\texttt{uStepScale} & 全局步长缩放 $\sigma_h$，与 $1/tol$ 近似等效（见 \S\ref{sec:perf}） \\
\texttt{btan < 1e-5} & 径向退化光线（$\sin\psi\to0$）走解析分支，不进入积分循环 \\
\texttt{int K = nearDisk ? 5 : 1} & 近盘段的一维子步数，用于把体积发射积到亚像素精度 \\
\texttt{dtau = min(\ldots, 6.0)} & 单段光学厚度上限，防止 $\exp(-\tau)$ 下溢导致透明度突跳 \\
\texttt{trans > 0.002} & 透射率截断：不透明后不再采样盘，节省算力 \\
\texttt{u < 1e-6} & 数值上「跑到无穷远」的兜底退出 \\
\bottomrule
\end{tabular}

\end{table}


\subsection{终止判据与光线分类}\label{sec:termination}

积分循环在五种情形下退出，按出现顺序为：

\begin{enumerate}\itemsep2pt
\item \textbf{捕获}：$u>0.5$，即 $r<\rs$。此时光线已经越过视界，
  序参量不再是「是否逃逸」而是「不再有逃逸的可能」，直接输出纯黑。
\item \textbf{逃逸}：$u<u_{\min}$ 且 $w<0$。$u_{\min}=1/\max(R_{\rm esc},1.6r_0)$，
  默认 $R_{\rm esc}=140\M$。要求 $w<0$（径向外行）是为了避免把
  尚未到达最近点、但恰好位于 $R_{\rm esc}$ 之内的外行段误判为逃逸；
  要求 $R_{\rm esc}\ge1.6r_0$ 则是为了在相机贴近视界时仍留出足够
  的「远场」余量，使 \S\ref{sec:v5} 的出射方向误差可控。
\item \textbf{数值无穷远}：$u<10^{-6}$（$r>10^6\M$），兜底退出。
\item \textbf{步数上限}：迭代计数达到运行期上限 $n_{\max}$
  （默认 $300$，界面可调 $60$--$900$）。
\item \textbf{编译期硬上限}：循环体写死 $4096$ 次。
  这一层保护与用户参数无关，用于防止在极端参数
  （$tol$ 极小、$\sigma_h$ 极小、$n_{\max}$ 极大）下
  单条病态光线挂死 GPU 触发 TDR。
\end{enumerate}

命中第 4 或第 5 条的光线按「未逃逸」处理：若此时 $u<0.5$，
它仍会把剩余透射率乘上背景星场（式~\eqref{eq:escapedir}），
因此截断表现为星场亮度与方向的微小误差，而不是伪影。
\S\ref{sec:perf} 表明：在默认参数下绝大多数像素远未逼近 $n_{\max}$，
$n_{\max}$ 从 $300$ 提高到 $1200$ 几乎不改变耗时（表~\ref{tab:perfsteps}），
这反过来印证了步长控制而非步数上限才是真正决定成本的因素。

\subsection{盘的体积辐射转移积分}\label{sec:rtnum}

盘不是几何薄面，而是光学厚度有限的体积发射体：密度在竖直方向为高斯
$\rho\propto\exp(-\tfrac12 y^2/H^2)$，$H=\mathcal H r$，
$\mathcal H$ 即界面上的 \texttt{diskThick}（默认 $0.075$）。
沿一条几何步长对应的线段 $[\vec p_1,\vec p_2]$ 上，
模拟器按前向（front-to-back）格式累积：
\begin{equation}
\dd\tau=\min\!\bigl(\kappa\rho\,\overline{w}_y\,\Delta \ell,\;6\bigr),
\qquad
\vec C \mathrel{+}= T\,(1-e^{-\dd\tau})\,\vec S,\qquad
T \mathrel{\times}= e^{-\dd\tau},
\label{eq:front2back}
\end{equation}
其中 $\Delta\ell=|\vec p_2-\vec p_1|$ 为线段长度，$\vec S$ 为
式~\eqref{eq:Tobs}--\eqref{eq:I4} 所定义的发射函数，$\overline{w}_y$ 为该段上竖直
高斯剖面的平均
\begin{equation}
\overline{w}_y=
\frac{1}{\Delta y}\int_{y_1}^{y_2}
\exp\!\Bigl(-\frac{y^2}{2H^2}\Bigr)\dd y
=\frac{1}{\Delta y}\sqrt{\frac{\pi}{2}}\,H\,
\Bigl[\mathrm{erf}\bigl(\tfrac{y_2}{H\sqrt2}\bigr)
-\mathrm{erf}\bigl(\tfrac{y_1}{H\sqrt2}\bigr)\Bigr],
\label{eq:avgw}
\end{equation}
$\Delta y\to0$ 时退化为 $\exp(-\tfrac12y^2/H^2)$。式~\eqref{eq:avgw}
把「穿过高斯壳层」的柱密度解析地积掉，因此即使一个几何步长横跨
整个薄盘，光学厚度仍然是准确的，不需要把 $h$ 缩到 $H$ 以下。
这一点对效率至关重要：默认 $H/R=0.075$、$r\sim10\M$ 时 $H\sim0.75\M$，
而远场步长对应的弧长可达几个 $M$。

$\mathrm{erf}$ 用 Abramowitz--Stegun 7.1.26 的有理逼近
（绝对误差 $<1.5\times10^{-7}$），在单精度下与精确值不可区分。
近盘区域（线段两端点中较近者落在 $7H$ 内、且半径位于
$[0.85\,\bin,1.3\,\bout]$）启用 $K=5$ 的等分子步，
把一段几何步长切成五份分别求值，使发射的空间变化被正确采样；
盘外区域则取 $K=1$，不做无用的细分。
当透射率 $T<0.002$ 时停止采样盘——光学厚之后继续积分只会
增加成本而不再改变像素值。

\subsection{守恒残差与收敛判据}\label{sec:residual}

式~\eqref{eq:orbit1} 的第一积分在解析解中严格守恒。数值解中定义残差
\begin{equation}
\mathcal R(\varphi)
=b^{-1}\Bigl[\bigl(\tfrac{\dd u}{\dd\varphi}\bigr)^{2}
+u^{2}(1-2\M u)\Bigr]^{1/2}-1,
\label{eq:resid}
\end{equation}
它以相对形式度量「这条数值曲线偏离真实零测地线多远」，
且对尺度不敏感。$\mathcal R$ 是后文最主要的内部一致性指标：
V2 在近临界轨道上报告它的最大值与平均值（表~\ref{tab:v2}），
图~\ref{fig:05} 则给出 $\max|\mathcal R|$ 随步数的收敛曲线。

残差有两条下界，读者在解读时需要注意：其一，着色器使用单精度
浮点（\texttt{highp float}，24 位尾数，机器精度 $\approx6\times10^{-8}$），
$\mathcal R$ 不可能长期低于该量级；其二，近临界轨道绕光子球多圈后
（V2 的测试光线绕行 $15$\,rad 以上）局部 Lyapunov 放大使得
初值的舍入噪声被指数放大，残差会在某一步「抬起」到
$10^{-6}$ 量级并保持。这正是表~\ref{tab:v2} 中残差随 $tol$ 收紧而下降
约两个数量级、随后在 $tol=3\times10^{-4}$ 触底的原因；
它反映的是浮点精度的极限，而不是积分格式的失误。

最后指出一处与图像质量直接相关的数值性质：式~\eqref{eq:hstep} 的
下界 $0.002$ 使得近光子球光线的步长不再由 $tol$ 决定，
而是固定为 $0.002\,\sigma_h$。这类光线在屏幕上恰好是阴影边缘
附近的少数像素，因此「阴影边缘的锐利度」几乎与 $tol$ 无关，
而「阴影边缘的位置」由式~\eqref{eq:bfromic} 的守恒律保证。
两者合起来解释了图~\ref{fig:17}(b) 中阴影半径随 $tol$ 迅速收敛而
图~\ref{fig:22}(c) 中总耗时仍随 $tol$ 显著变化的表观矛盾。
